\uuid{Qc6c}
\titre{Primitives de fonctions rationnelles }

\niveau{} 				%L1, L2, L3, MPSI, MP, PCSI, PC, PSI...
\module{ Intégration } 	%Analyse, Algèbre...
\chapitre{Méthode d'intégration}   			%Continuité, Groupes, Fonctions de plusieurs variables...
\sousChapitre{}			%Optimisation, Diagonalisation d'une matrice, Calcul de dérivées partielles...

\theme{}				%Fonction de répartition, Division euclidienne de polynômes, ...
\auteur{Quentin Liard}
\datecreate{ 2026-02-11}
\organisation{AMSCC}			%AMSCC, Exo7, ...
\difficulte{2}			%1, 2, 3, 4 ou 5

\contenu{

\texte{ 

}

\begin{enumerate}
\item   \question{Calculer $I_1=\displaystyle\int \frac{1}{x^2+1}\,dx$.}
\indication{}
\reponse{$\displaystyle I_1=\arctan x + C$.}

\item   \question{Montrer une relation entre $I_2=\displaystyle\int \frac{1}{(x^2+1)^2}\,dx$ et $I_1$ en utilisant (par exemple) l’identité $\dfrac{d}{dx}\!\left(\dfrac{x}{1+x^2}\right)=\dfrac{1-x^2}{(1+x^2)^2}$.}
\indication{}
\reponse{On a $\displaystyle \frac{1}{(1+x^2)^2}=\frac12\left(\frac{1}{1+x^2}+\frac{1-x^2}{(1+x^2)^2}\right)=\frac12\left(\frac{1}{1+x^2}+\frac{d}{dx}\left(\frac{x}{1+x^2}\right)\right)$. En intégrant : $I_2=\frac12 I_1+\frac12\frac{x}{1+x^2}+C$.}

\item   \question{En déduire une expression explicite de $I_2$.}
\indication{}
\reponse{Comme $I_1=\arctan x + C$, on obtient une primitive : $\displaystyle I_2=\frac12\arctan x+\frac12\frac{x}{1+x^2}+C$.}

\end{enumerate}

}